\section{Introduction}
\label{sec:intro}

\begin{figure}[t]
  \centering
  % Figure 1: the extension around the partitioned CSR, framed by the engine hooks it uses.
% Colours follow paper/figures: operator blue, index orange, engine grey. Coordinates in cm; text width 17.8 cm.
\begin{tikzpicture}[
  font=\scriptsize,
  >={Stealth[length=4pt,width=3.5pt]},
  box/.style={draw, rounded corners=1.5pt, align=center, inner sep=2.5pt, minimum height=5mm, anchor=north west},
  engine/.style={box, draw=engine!80!black, fill=engine!15, text=black!80},
  hook/.style={box, draw=engine!80!black, fill=white, font=\scriptsize\itshape, inner sep=2pt, minimum height=4mm, anchor=north},
  op/.style={box, draw=opblue, fill=opblue!10},
  ix/.style={box, draw=ixorange, fill=ixorange!10},
  sql/.style={box, draw=black!50, fill=white, font=\tiny\ttfamily, align=left},
  flow/.style={->, thick, black!70},
  opflow/.style={->, thick, opblue},
  ixflow/.style={->, thick, ixorange},
  hooked/.style={->, densely dashed, engine!80!black},
  lbl/.style={font=\tiny, text=black!60, inner sep=1pt},
]
% ---------------------------------------------------------------- engine band (top)
\node[engine, minimum width=17.8cm, align=left] (band) at (0,0)
  {\textbf{DuckDB} \quad parser \;$\cdot$\; binder \;$\cdot$\; optimizer \;$\cdot$\; pipelines and scheduler \;$\cdot$\; MVCC storage \;$\cdot$\; index manager \;$\cdot$\; checkpoint and WAL};
\node[hook] (h-grammar) at (2.3,-0.75)  {grammar extension};
\node[hook] (h-opt)     at (7.6,-0.75)  {optimizer extension};
\node[hook] (h-physop)  at (12.9,-0.75) {physical operator};
\node[hook] (h-index)   at (17.0,-0.75) {index type};
\foreach \h in {h-grammar,h-opt,h-physop,h-index} {\draw[engine!80!black] (\h.north) -- (\h.north |- band.south);}

% ---------------------------------------------------------------- row 1: query rewrite
\node[sql, text width=4.6cm] (match) at (0,-1.55)
  {MATCH p = ANY SHORTEST\\ \ (a:Person)-[k:Knows]->\{1,8\}(b:Person)\\ \ WHERE a.id = 7\\ COLUMNS (b.id, path\_length(p))};
\node[sql, text width=4.2cm] (marker) at (5.5,-1.55)
  {SELECT b.id, iterativelength\\ \ \ \ (a.id, b.id, edges, V, E, key, 8)\\ FROM Person a, Person b, edges\\ WHERE a.id = 7};
\node[op, text width=4.7cm, minimum height=11mm] (operator) at (10.6,-1.55)
  {\textbf{path-finding operator}\\ sink: pairs $(a,b)$, edge endpoints\\ source: one length or path per pair};
\draw[flow] (match) -- node[lbl, above] {rewrite} (marker);
\draw[flow] (marker) -- node[lbl, above] {replace} node[lbl, below] {marker} (operator);
\draw[hooked] (h-grammar.south) -- (h-grammar.south |- match.north);
\draw[hooked] (h-opt.south) -- (h-opt.south |- marker.north);
\draw[hooked] (h-physop.south) -- (h-physop.south |- operator.north);

% ---------------------------------------------------------------- row 2: build per query -> partitioned CSR <- index
\node[op, text width=4.9cm] (build) at (0,-3.55)
  {\textbf{build per query} (\S\ref{sec:csr})\\ edge table $\bowtie$ vertices $\rightarrow$ endpoints\\ radix partition by destination range\\ sort by source, emit sparse rows};
\node[box, draw=black!60, fill=white, inner sep=2pt] (csr) at (6.0,-3.55) {%
  \begin{tikzpicture}[font=\tiny, inner sep=0pt]
    \node[anchor=north west, text=black!70, font=\scriptsize] (title) at (0,0) {\textbf{partitioned CSR}: one partition per destination range};
    \foreach \i/\t in {0/0,1/1,2/2,3/$\cdots$,4/$p$} {
      \node[draw=black!60, fill=black!4, minimum width=5.4mm, minimum height=4mm, anchor=north west, font=\scriptsize] (p\i) at ($(title.south west)+(\i*5.7mm,-1mm)$) {\t};
    }
    \node[draw=black!60, fill=white, anchor=north west, align=left] (detail) at ($(p0.south west)+(0,-1.4mm)$)
      {sparse rows: source id, offsets, bucket index\\ edges: 16-bit destination (+ row id for paths)};
    \draw[black!60] (p1.south) -- (p1.south |- detail.north);
  \end{tikzpicture}};
\node[ix, text width=4.9cm] (index) at (12.6,-3.55)
  {\textbf{CSR index} (\S\ref{sec:index})\\ immutable epoch CSR + dictionary\\ change log stamped per commit group\\ snapshot = epoch + overlay for txn $T$};
\draw[opflow] (build.east) -- node[lbl, above] {per query} (build.east -| csr.west);
\draw[ixflow] (index.west) -- node[lbl, above] {snapshot} (index.west -| csr.east);
\draw[hooked] (h-index.south) -- (h-index.south |- index.north);
\draw[opflow] (operator.south) -- ++(0,-2.5mm) -| ($(csr.north east)+(-4mm,0)$);

% ---------------------------------------------------------------- row 3: traversals and result
\node[op, text width=4.3cm] (ms) at (0,-5.75)
  {\textbf{multi-source traversal} (\S\ref{sec:traversals})\\ 512 pairs per batch in bit lanes\\ one sweep over all partitions per level};
\node[op, text width=4.6cm] (ss) at (5.0,-5.75)
  {\textbf{single-source traversal} (\S\ref{sec:traversals})\\ frontier list, bucket lookups, parents\\ stops at the depth bound or last destination};
\node[op, text width=3.9cm] (phases) at (10.3,-5.75)
  {\textbf{phases as events} (\S\ref{sec:phases})\\ one event per BFS level\\ no task waits for a task};
\node[sql, text width=2.9cm] (result) at (14.9,-5.75)
  {b.id  len  path\\ 42 \ \ \ 3 \ \ [7,e1,19,e2,42]\\ 77 \ \ \ 5 \ \ ...};
\draw[opflow] (csr.south) -- ++(0,-2.5mm) -| ($(ss.north)+(0,0)$);
\draw[opflow] (ms.east) -- (ss.west);
\draw[opflow] (ss.east) -- (phases.west);
\draw[flow] (phases.east) -- node[lbl, above] {filter} node[lbl, below] {on \{1,8\}} (result.west);
\end{tikzpicture}

  \caption{Path finding in \duckpgq. A \texttt{MATCH} path pattern is rewritten to SQL with a marker function that an
  optimizer extension replaces by the path-finding operator (Section~\ref{sec:planner}). The operator traverses a
  partitioned CSR that is either built inside the query from the edge table (Section~\ref{sec:csr}) or read as a
  transaction-consistent snapshot of a CSR index (Section~\ref{sec:index}); both have the same layout. Two
  traversals, multi-source lanes for many pairs and a frontier search for few, run as chains of per-phase events
  (Sections~\ref{sec:traversals} and~\ref{sec:phases}). Everything in colour is the extension; the grey band is
  \duckdb, with the four of its extension interfaces that carry the query (Section~\ref{sec:strategy} lists all
  of them).}
  \label{fig:overview}
\end{figure}

% Draft 22 Oct 2026 in the style of paper/STYLE.md. Every number is provisional until the frozen-commit rerun.

SQL:2023 added graph pattern matching to SQL, and with it a query type relational systems have never executed
well, the shortest path. Its \sqlpgq part lets a query declare a property graph over existing tables and match
patterns against it, including variable-length and shortest-path patterns~\cite{deutsch2022gql}. A user whose data
lives in a relational system has three options for such queries. Copy the data into a graph database and keep the
two in sync. Write the path pattern as a recursive join and accept its performance. Or have the system answer it
with graph storage and operators of its own, which so far has meant waiting for its developers to build them. This
paper shows that a loadable extension can add them.

\duckpgq~\cite{tenwolde2023duckpgq} was our first attempt. It is an extension of \duckdb that parses \sqlpgq and
rewrites it into relational SQL. We argued there that a graph system should build on everything a state-of-the-art
analytical system already has, columnar storage, vectorized execution, morsel-driven parallelism and a cost-based
optimizer, and add two things, a compact graph representation and a bulk path-finding algorithm. We implemented both
as scalar user-defined functions (UDFs). One set of UDFs builds a compressed sparse row (CSR) representation on the
fly while the query streams the edge table past them; another runs multi-source breadth-first search
(\msbfs)~\cite{then2014msbfs} over that CSR, 512 source-destination pairs per invocation. UDFs were the right choice
for a first system. They are as fast as built-in functions, they inherit \duckdb's parallelism, and they kept the
extension small enough to follow \duckdb's release cadence.

They also have a ceiling, which the CIDR paper named: ``morsel-driven parallelism is tuple-driven and without tuples
there is not enough parallelism''~\cite{tenwolde2023duckpgq}. A UDF receives work when the pipeline hands it a vector
of tuples. A single shortest-path query hands it one tuple, so one core runs a BFS over the whole graph while the
other fifteen idle, and with 16K pairs the same code scaled by a factor six on 16 threads. The CSR is a second problem
the UDF design could not solve. It is global state built by ``the first thread to touch it'', rebuilt for every
query, invisible to the optimizer, and unrelated to the transaction that reads it, so a concurrent insert into the
edge table is in the CSR or not depending on timing. The paper left both as a research challenge, ``to reconcile this
elegantly with any query-pipeline-driven parallelization method''.

This paper takes up that challenge under the constraint that made \duckpgq viable, that everything stays in a
loadable extension. Our thesis is that the interfaces \duckdb already offers to extensions are enough to give it
graph-native path finding, including a graph index that is consistent with its multi-version concurrency control
(MVCC), and that the places where they fall short are few and can be named.

The first half of the answer is a \emph{path-finding operator}. It is a physical operator with its own inputs and its
own schedule, so its parallelism comes from the graph and not from the tuples. The edges are partitioned by
destination range, the partitions are built and traversed in parallel, and the traversal proceeds level by level as a
chain of events in which no task blocks its thread waiting for another. One traversal serves many pairs at once, the multi-source lanes of the
CIDR design; a second, a frontier search per source, serves the single pair that the lanes handle badly. The planner
reaches the operator through an optimizer rule and sees an ordinary operator with a cardinality.

The second half is a \emph{CSR index}. \texttt{CREATE INDEX ... USING CSR (src, dst)} keeps the graph representation
across queries, restarts and checkpoints, maintained through \duckdb's index interface like any other index. The
difficulty is consistency. A transaction must traverse exactly the edges it is allowed to see, and the interface
tells an extension which rows changed but nothing about who may see them. The index stores immutable epochs and a
change log whose entries \duckdb's own row-version check classifies, once per commit instead of once per edge. Every
transaction gets a snapshot of its own transaction's view, a writer spends a constant amount of work per edge, and a
background fold turns the log into the next epoch. Section~\ref{sec:index} gives the design and its correctness
argument.

We evaluate on LDBC SNB scale factors 1 to 300 and the Graphalytics data sets against Kuzu, Neo4j, Memgraph and
FalkorDB, read-only and under LDBC update streams with concurrent transactions. \todo{FINAL: pick three numbers, one per
claim (single pair over the index, many pairs vs the native systems, update cost per commit). Provisional,
m7a.8xlarge, 32 threads: at SF300 a single-source query takes 0.30\,s when the CSR is built per query and 0.031\,s
over the index; the index loads from disk in 0.09\,s; under a stream of 2000 edge updates every 200\,ms with 16
concurrent transactions, the index answers single-source queries in 72\,ms and the build-per-query route in 84\,ms, both
with results identical to Kuzu at every step; Kuzu needs 16\,ms per query but 342\,ms per update batch where \duckdb
needs 9. Replace with frozen-commit numbers.}

Our contributions are:
\begin{enumerate}
  \item A path-finding operator for a vectorized, morsel-driven system whose parallelism comes from a partitioned
        CSR and per-level events instead of from input tuples, with multi-source and single-source traversals and a
        rule that selects between them (Section~\ref{sec:operator}).
  \item An MVCC-consistent CSR index built behind a general-purpose extension index interface: epochs, a commit-group
        change log, per-transaction snapshots, folds, checkpoints and recovery, with a correctness argument
        (Section~\ref{sec:index}).
  \item The planner integration from \sqlpgq \texttt{MATCH} to operator and index without changes to \duckdb, and the
        list of what a system would have to expose to make such an extension simpler (Section~\ref{sec:strategy}).
  \item An evaluation against four native graph systems, read-only and under updates, showing that \duckdb
        with the index matches them on point and single-source queries while keeping every update in the
        same transaction as the tables (Section~\ref{sec:eval}).
\end{enumerate}
\duckpgq is open source under the MIT licence\footnote{\url{https://github.com/cwida/duckpgq-extension}}; the
operator and the index are part of it, and the benchmark harness and result archive are released with this paper.
